\documentclass{beamer}
\usetheme{Madrid}
\usepackage{calc}
\usepackage{graphicx}
\usepackage{amsmath}
\usepackage{amssymb}
\usepackage{relsize}
\usepackage{multirow}
\usepackage{rotating}
\usepackage{bm}
\usepackage{tikz}
\usetikzlibrary{tikzmark,positioning,calc}

\usepackage[natbib=true,style=authoryear,backend=bibtex,useprefix=true]{biblatex}
\addbibresource{refs.bib}

\graphicspath{{images/}{../images/}}
%\definecolor{lightblue}{cmyk}{0.83,0.24,0,0.12}
\definecolor{lightblue}{rgb}{0.145,0.6666,1}

% Draw a video
\newlength{\FSZ}
\newcommand{\drawvideo}[3]{% [0 0.25 0.5 0.75 1 1.25 1.5]
	\noindent\pgfmathsetlength{\FSZ}{\linewidth/#2}
	\begin{tikzpicture}[outer sep=0pt,inner sep=0pt,x=\FSZ,y=\FSZ]
	\draw[color=lightblue!50!black] (0,0) node[outer sep=0pt,inner sep=0pt,text width=\linewidth,minimum height=0] (video) {\noindent#3};
	\path [fill=lightblue!50!black,line width=0pt] 
	(video.north west) rectangle ([yshift=\FSZ] video.north east) 
	\foreach \x in {1,2,...,#2} {
		{[rounded corners=0.6] ($(video.north west)+(-0.7,0.8)+(\x,0)$) rectangle +(0.4,-0.6)}
	}
	;
	\path [fill=lightblue!50!black,line width=0pt] 
	([yshift=-1\FSZ] video.south west) rectangle (video.south east) 
	\foreach \x in {1,2,...,#2} {
		{[rounded corners=0.6] ($(video.south west)+(-0.7,-0.2)+(\x,0)$) rectangle +(0.4,-0.6)}
	}
	;
	\foreach \x in {1,...,#1} {
		\draw[color=lightblue!50!black] ([xshift=\x\linewidth/#1] video.north west) -- ([xshift=\x\linewidth/#1] video.south west);
	}
	\foreach \x in {0,#1} {
		\draw[color=lightblue!50!black] ([xshift=\x\linewidth/#1,yshift=1\FSZ] video.north west) -- ([xshift=\x\linewidth/#1,yshift=-1\FSZ] video.south west);
	}
	\end{tikzpicture}
}



\title[Test of pRF Estimates] %optional
{Test of Goodness for \\Population Receptive Field Estimates}

\subtitle{simulation study}

\author[Ashrafnejad, Arash] % (optional, for multiple authors)
{A.~Ashrafnejad\inst{1} \and H.~Mehrzadfar\inst{2} \and H.~Boyaci\inst{2} \inst{3}}

\institute[CBVG] % (optional)
{
	\inst{1}%
	Department of Electrical and Electronics Engineering
	\and
	\inst{2}%
	Neuroscience Graduate Program
	\and
	\inst{3}%
	Department of Psychology\\
	
	Bilkent University
}

\date{\today}


\begin{document}
	
	\begin{frame}
	\titlepage
	\end{frame}


	\section{Data Generation}
	\begin{frame}{Stimulation}
	  \noindent We first used the exact drifting bar stimululation in \cite{Dumoulin2008}:
	
%	\drawvideo{24}{72}{%
%		\includegraphics[width=0.0417\linewidth]{{stim_cont_0.00_0}.png}%
%		\includegraphics[width=0.0417\linewidth]{{stim_cont_0.00_1}.png}%
%		\includegraphics[width=0.0417\linewidth]{{stim_cont_0.00_2}.png}%
%		\includegraphics[width=0.0417\linewidth]{{stim_cont_0.00_3}.png}%
%		\includegraphics[width=0.0417\linewidth]{{stim_cont_0.00_4}.png}%
%		\includegraphics[width=0.0417\linewidth]{{stim_cont_0.00_5}.png}%
%		\includegraphics[width=0.0417\linewidth]{{stim_cont_0.00_6}.png}%
%		\includegraphics[width=0.0417\linewidth]{{stim_cont_0.00_7}.png}%
%		\includegraphics[width=0.0417\linewidth]{{stim_cont_0.00_8}.png}%
%		\includegraphics[width=0.0417\linewidth]{{stim_cont_0.00_9}.png}%
%		\includegraphics[width=0.0417\linewidth]{{stim_cont_0.00_10}.png}%
%		\includegraphics[width=0.0417\linewidth]{{stim_cont_0.00_11}.png}%
%		\includegraphics[width=0.0417\linewidth]{{stim_cont_0.00_12}.png}%
%		\includegraphics[width=0.0417\linewidth]{{stim_cont_0.00_13}.png}%
%		\includegraphics[width=0.0417\linewidth]{{stim_cont_0.00_14}.png}%
%		\includegraphics[width=0.0417\linewidth]{{stim_cont_0.00_15}.png}%
%		\includegraphics[width=0.0417\linewidth]{{stim_cont_0.00_16}.png}%
%		\includegraphics[width=0.0417\linewidth]{{stim_cont_0.00_17}.png}%
%		\includegraphics[width=0.0417\linewidth]{{stim_cont_0.00_18}.png}%
%		\includegraphics[width=0.0417\linewidth]{{stim_cont_0.00_19}.png}%
%		\includegraphics[width=0.0417\linewidth]{{stim_cont_0.00_20}.png}%
%		\includegraphics[width=0.0417\linewidth]{{stim_cont_0.00_21}.png}%
%		\includegraphics[width=0.0417\linewidth]{{stim_cont_0.00_22}.png}%
%		\includegraphics[width=0.0417\linewidth]{{stim_cont_0.00_23}.png}%
%	}
	\end{frame}

	\begin{frame}{fMRI Signal Modelling}
	
	\begin{exampleblock}{BOLD response model}
		\[
		B(t) =  \tikzmark{first}\mathcal{H}(\quad \tikzmark{second}r(t, \Theta=\tikzmark{third}\theta)) + \tikzmark{forth}e(t)
		\]
		\begin{tikzpicture}[
		remember picture,
		overlay,
		expl/.style={draw=orange,fill=orange!30,rounded corners,text width=3cm},
		arrow/.style={red!80!black,ultra thick,->,>=latex}
		]
		\node<2->[expl] 
		(firstN) 
		at (4,-2cm)
		{Linear or Non-linear HRF function};
		
		\node<3->[expl] 
		(secondN) 
		at (4,3.5cm)
		{assumed pRF response (neuronal response)};
		
		\node<4->[expl] 
		(thirdN) 
		at (8,-2cm)
		{assumed pRF parameters; $x_0,y_0,\sigma$};
		
		\node<5->[expl] 
		(forthN) 
		at (8,3.5cm)
		{Gaussian noise due to large population of neurons};
		
		\draw<2->[arrow]
		(firstN) to[out=90,in=-90] ([yshift=0.5ex]{pic cs:first});  
		
		\draw<3->[arrow]
		(secondN.south) to[out=-90,in=90] ([yshift=0.5ex]{pic cs:second});  
		
		\draw<4->[arrow]
		(thirdN.north) to[out=90,in=-90] ([yshift=0.5ex]{pic cs:third});  
		
		\draw<5->[arrow]
		(forthN.south) to[out=-90,in=90] ([yshift=0.5ex]{pic cs:forth});  
		\end{tikzpicture}
	\end{exampleblock}
	\end{frame}
	
	\begin{frame}{Neuronal Response Modelling}
	\begin{itemize}
		\item 
		Neurons within a small region of visual cortex respond to stimuli within a restricted region of the visual field.
		\item
		The population response of such neurons can $\boldsymbol{not}$ be modeled using a model that sums contrast linearly across the visual field. hence, Compressive spatial summation (CSS) model is used \cite{Kay2013},
	\end{itemize}
	\begin{exampleblock}{pRF response model}
	\[
	r(t, \Theta=\theta) = (\sum_{x, y} \tikzmark{first2}s(x, y, t) \tikzmark{second2} g(x,y, \theta))^ n \tikzmark{third2}
	\]
	\begin{tikzpicture}[
	remember picture,
	overlay,
	expl/.style={draw=orange,fill=orange!30,rounded corners,text width=3cm},
	arrow/.style={red!80!black,ultra thick,->,>=latex}
	]
	\node<2->[expl] 
	(firstN2) 
	at (2,-1cm)
	{stimulus function to be presented};
	
	\node<3->[expl] 
	(secondN2) 
	at (6,-1cm)
	{assumed pRF model};
	
	\node<4->[expl] 
	(thirdN2) 
	at (10,-1cm)
	{spatial linearity factor from 0 to 1};

	
	\draw<2->[arrow]
	(firstN2) to[out=90,in=-90] ([yshift=0.5ex]{pic cs:first2});  
	
	\draw<3->[arrow]
	(secondN2.north) to[out=90,in=-90] ([yshift=0.5ex]{pic cs:second2});  
	
	\draw<4->[arrow]
	(thirdN2.north) to[out=90,in=-90] ([yshift=0.5ex]{pic cs:third2});  

	\end{tikzpicture}
	\end{exampleblock}
	\end{frame}

	\begin{frame}{pRF Modelling}
		\includegraphics[height=7.0em]{2d.png}
		$\sigma = \frac{1}{2} \ln(e + \sqrt{x_0^2 + 2y_0^2})$
			\begin{exampleblock}{pRF response model}
			\[
			g(x, y, \Theta=\theta) = e^{-\frac{(x-\tikzmark{first3}x_0)^2+(y-\tikzmark{second3}y_0)^2}{2\sigma^2\tikzmark{third3}}}
			\]
			\begin{tikzpicture}[
			remember picture,
			overlay,
			expl/.style={draw=orange,fill=orange!30,rounded corners,text width=3cm},
			arrow/.style={red!80!black,ultra thick,->,>=latex}
			]
			\node<2->[expl] 
			(firstN3) 
			at (6,-1cm)
			{assumed pRF center for the voxel};
			
			
			\node<4->[expl] 
			(thirdN3) 
			at (10,4cm)
			{pRF size assumed to be a log-polar function of pRF center};
			
			
			\draw<2->[arrow]
			(firstN3) to[out=90,in=-90] ([yshift=0.5ex]{pic cs:first3});  
			
			\draw<3->[arrow]
			(firstN3.north) to[out=90,in=-90] ([yshift=0.5ex]{pic cs:second3});  
			
			\draw<4->[arrow]
			(thirdN3.south) to[out=-90,in=0] ([yshift=0.5ex]{pic cs:third3});  
			
			\end{tikzpicture}
		\end{exampleblock}	\end{frame}


	\begin{frame}[t,allowframebreaks]
	\frametitle{References}
	\printbibliography
	\end{frame}


\end{document}